% ====================================================================
\chapter{Summation Notation}\label{ch:sum}
% ====================================================================

\begin{chapterintro}
Totals are central to economics: total revenue over twelve months, total consumption over a lifetime, the sum
of the incomes of everyone in a country. Writing out such sums term by term is impractical, and the summation
operator $\sum$ provides a compact and precise notation for them. This chapter defines the notation, fixes the
conventions for unusual cases, and establishes the three properties that allow sums to be manipulated.
\end{chapterintro}

\begin{prerequisites}
The properties of addition and multiplication A1--A4, M1--M5 (\cref{ch:numbers}); the natural numbers and integers
(\cref{sec:number-systems}).
\end{prerequisites}

\section{The summation operator}\label{sec:sigma}

We can add many numbers, such as $3+5+2+9+5$. When the numbers are labelled by an index that runs through
consecutive integers, there is a shorthand.

\begin{definition}[Summation operator][def:sigma]
Let $a_k,a_{k+1},\dots,a_n$ be numbers, indexed by an integer $j$ that takes the consecutive values
$k,k+1,\dots$ up to $n$. Their sum is written
\[
\sum_{j=k}^{n}a_j=a_k+a_{k+1}+\dots+a_{n-1}+a_n.
\]
We read this as follows: as the index $j$ changes from $k$ to $n$, we sum up all the corresponding numbers $a_j$.
In running text the same sum is often written $\sum_{j=k}^n a_j$. \src{L1, slide 27}
\end{definition}

The integer $j$ is the \term{index of summation}, $k$ is the \term{lower limit} and $n$ the \term{upper limit}.
The index is a \emph{dummy variable}: it has no meaning outside the sum, and any other letter can be used without
changing the value, $\sum_{j=1}^4 j^2=\sum_{i=1}^4 i^2=30$. The limits $k$ and $n$ may be any integers, including
negative ones and zero.

\begin{example}[Evaluating sums][ex:sums]
Evaluate, for $a_j=j$ and $b_j=j^2$,
\[
\sum_{j=3}^{7}a_j,\qquad\sum_{j=-2}^{3}a_j,\qquad\sum_{j=3}^{7}b_j,\qquad\sum_{j=-2}^{3}b_j.
\]
\src{L1, slide 28}
\solution
Write out each term and add:
\begin{align*}
\sum_{j=3}^{7}j&=3+4+5+6+7=25,\\
\sum_{j=-2}^{3}j&=-2-1+0+1+2+3=3,\\
\sum_{j=3}^{7}j^2&=3^2+4^2+5^2+6^2+7^2=9+16+25+36+49=135,\\
\sum_{j=-2}^{3}j^2&=(-2)^2+(-1)^2+0^2+1^2+2^2+3^2=4+1+0+1+4+9=19.
\end{align*}
Note that the sums starting at $j=-2$ include the terms for $j=-2,-1$ and $0$.
\end{example}

\section{Counting terms, one-term sums and empty sums}\label{sec:sigma-conventions}

\begin{core}[The number of terms]
The sum $\sum_{j=k}^n a_j$ (with $n\ge k$) has $n-k+1$ terms, because both the lower and the upper limit are
included. For example, $\sum_{j=3}^7$ has $7-3+1=5$ terms ($j=3,4,5,6,7$), not $4$.
\end{core}

Two boundary cases need to be fixed.

\begin{definition}[One term and zero terms][def:sigma-boundary]
\begin{itemize}
  \item When $n=k$, the sum consists of one term, with index $j=n=k$: \ $\displaystyle\sum_{j=k}^{k}a_j=a_k$.
  \item By analogy, when $n=k-1$, we define the value of the summation operator as a sum of zero terms, which we
  take to be zero: \ $\displaystyle\sum_{j=k}^{k-1}a_j=0$.
\end{itemize}
\src{L1, slide 29}
\end{definition}

The first is a consequence of the definition; the second is a \emph{convention}, a choice made so that the notation
behaves well. Why zero? Adding a sum of no numbers should not change a total, and zero is the number whose addition
changes nothing (A3). \Cref{prop:empty-forced} below shows that zero is in fact the only value consistent with the
splitting property P3.

\begin{keyformula}[Sum of a constant][kf:constant-sum]{\sum_{j=k}^{n}a=\underbrace{a+a+\dots+a}_{n-k+1\ \text{terms}}=(n-k+1)\,a}
\fsymbols $a$ is a fixed real number (the same for every $j$); $k\le n$ are integers; $n-k+1$ is the number of terms.
\fmeaning If every term equals $a$, the sum is the number of terms times $a$. \src{L1, slide 29}
\forigin Repeated addition is multiplication: the sum has $n-k+1$ terms, each equal to $a$, and
$a+a+\dots+a=(1+1+\dots+1)a$ by M5.
\fuse Whenever a sum contains a constant part, for example after splitting $\sum(4j+3)$ into $4\sum j+\sum 3$.
\fexample $\sum_{j=-5}^{9}4$: the count is $9-(-5)+1=15$, so the sum is $15\times4=60$.
\fmistakes Using $n-k$ instead of $n-k+1$ (which would give $56$ above); with a negative lower limit, computing
$9-5$ instead of $9-(-5)$.
\fchanges Raising the upper limit by one adds one more copy of $a$; for $n=k-1$ the formula gives $0\cdot a=0$, in
agreement with the empty-sum convention.
\fexam As a short computation (``evaluate $\sum_{j=-4}^{11}7$''), or as a step in finding a closed form.
\end{keyformula}

\section{Properties of the summation operator}\label{sec:sigma-props}

\begin{proposition}[Properties P1--P3][prop:sigma-props]
For numbers $a_j$, $b_j$ and a constant $c$:
\begin{align*}
\text{(P1)}\quad &\sum_{j=k}^{n}(a_j+b_j)=\sum_{j=k}^{n}a_j+\sum_{j=k}^{n}b_j,\\
\text{(P2)}\quad &\sum_{j=k}^{n}c\cdot a_j=c\cdot\sum_{j=k}^{n}a_j,\\
\text{(P3)}\quad &\sum_{j=k}^{n}a_j+\sum_{j=n+1}^{m}a_j=\sum_{j=k}^{m}a_j\qquad\text{for }k\le n\le m.
\end{align*}
P1 and P2 follow from the properties of the real numbers; P3 comes from the definition of the operator.
\src{L1, slide 30}
\end{proposition}
\emph{Proof.} (P1) The left side is $(a_k+b_k)+(a_{k+1}+b_{k+1})+\dots+(a_n+b_n)$. By commutativity and
associativity of addition (A1, A2), the terms may be reordered and regrouped as
$(a_k+\dots+a_n)+(b_k+\dots+b_n)$.

(P2) $ca_k+ca_{k+1}+\dots+ca_n=c(a_k+a_{k+1}+\dots+a_n)$ by the distributive law (M5, applied repeatedly).

(P3) The first sum adds the terms with indices $k,\dots,n$, the second those with indices $n+1,\dots,m$; together
they add each term with index from $k$ to $m$ exactly once. (When $n=m$ the second sum is empty and equals $0$.)
\hfill$\square$

The difference between the two kinds of justification is worth noticing. P1 and P2 are about the numbers being
added, and they hold because of how addition and multiplication behave. P3 is about the notation: splitting the
range of the index and re-joining it gives back the same list of terms.

\begin{proposition}[The empty sum must be zero][prop:empty-forced]
If P3 is to hold also when $n=k-1$, then $\sum_{j=k}^{k-1}a_j$ must equal $0$.
\end{proposition}
\emph{Proof.} Put $n=k-1$ in P3 with any $m\ge k$: $\sum_{j=k}^{k-1}a_j+\sum_{j=k}^{m}a_j=\sum_{j=k}^{m}a_j$.
Adding $-\sum_{j=k}^m a_j$ to both sides (A5, A4) gives $\sum_{j=k}^{k-1}a_j=0$. \hfill$\square$

\begin{warning}[What the properties do not say]
There is no rule for products or powers inside a sum:
\[
\sum_{j}a_jb_j\ne\Big(\sum_j a_j\Big)\Big(\sum_j b_j\Big),\qquad \sum_j a_j^2\ne\Big(\sum_j a_j\Big)^2\quad\text{in general.}
\]
For example, $\sum_{j=1}^2 j^2=1+4=5$, while $\big(\sum_{j=1}^2 j\big)^2=3^2=9$. Also, P2 allows only a factor $c$ that
does not depend on the index to be taken outside.
\end{warning}

\begin{external}[Shifting the index]
Replacing $j$ by $i+1$ throughout shifts the limits down by one:
$\sum_{j=k}^{n}a_j=\sum_{i=k-1}^{n-1}a_{i+1}$. Both sides list the same terms. For example
$\sum_{j=1}^{n}j=\sum_{i=0}^{n-1}(i+1)$. Index shifts are needed to combine sums with different limits.
\end{external}

\section{Computing sums with the properties}\label{sec:sigma-compute}

\begin{example}[A sum with a negative lower limit]
Evaluate $\sum_{j=-1}^{4}(j^2-j)$, showing the individual terms. \src{SG, mock Q4(a)}
\solution
For $j=-1,0,1,2,3,4$ the terms $j^2-j$ are
$(1+1)+(0-0)+(1-1)+(4-2)+(9-3)+(16-4)=2+0+0+2+6+12=22$.
\end{example}

\begin{example}[Two methods for one sum]
Compute $\sum_{j=2}^{6}(2j-1)$ directly and by the properties.
\solution
Directly: $3+5+7+9+11=35$. By the properties, with $\sum_{j=2}^6 j=2+3+4+5+6=20$ and $6-2+1=5$ terms:
\[
\sum_{j=2}^{6}(2j-1)\overset{\text{P1}}{=}\sum_{j=2}^{6}2j+\sum_{j=2}^{6}(-1)\overset{\text{P2}}{=}2\sum_{j=2}^{6}j+5\cdot(-1)=40-5=35.
\]
\end{example}

\begin{example}[A closed form][ex:closed-form]
Given that $\sum_{j=1}^{n}j=\frac{n(n+1)}{2}$ (proved in \cref{ch:ind}), find a closed-form expression for
$\sum_{j=1}^{n}(4j+3)$, naming each property used. \src{SG, mock Q4(b)}
\solution
\[
\sum_{j=1}^{n}(4j+3)\overset{\text{P1}}{=}\sum_{j=1}^{n}4j+\sum_{j=1}^{n}3\overset{\text{P2}}{=}4\sum_{j=1}^{n}j+\sum_{j=1}^{n}3
=4\cdot\frac{n(n+1)}{2}+(n-1+1)\cdot3=2n^2+2n+3n=2n^2+5n,
\]
using the constant-sum rule (\cref{kf:constant-sum}) for $\sum 3$. Check with $n=2$: $7+11=18$ and
$2\cdot4+5\cdot2=18$.
\end{example}

\begin{example}[Further evaluations]
Evaluate (a) $\sum_{j=-3}^{2}j^3$; (b) $\sum_{j=-4}^{11}7$; (c) $\sum_{j=7}^{6}a_j$.
\solution
(a) $(-27)+(-8)+(-1)+0+1+8=-27$. (The terms for $j=\pm1$ and $j=\pm2$ cancel in pairs.)
(b) $11-(-4)+1=16$ terms, so $16\times7=112$.
(c) The upper limit is one below the lower limit, so this is an empty sum, equal to $0$ by convention.
\end{example}

\begin{external}[Recalling $\sum_{j=1}^n j$ by pairing]
Write the sum forwards and backwards and add: each of the $n$ columns of $1+2+\dots+n$ and $n+(n-1)+\dots+1$ adds to
$n+1$, so twice the sum is $n(n+1)$. This is a way to \emph{recall} the formula; in the course it is \emph{proved} by
induction (\cref{ex:sum-n}).
\end{external}

\section{Sums in economics}\label{sec:sigma-econ}

\begin{example}[Total and average revenue]
A shop's revenue in month $t$ is $R_t$ for $t=1,\dots,12$. Write total annual revenue and average monthly revenue in
summation notation. If $R_t=1000+50t$, compute both.
\solution
Total revenue is $\sum_{t=1}^{12}R_t$ and the average is $\frac1{12}\sum_{t=1}^{12}R_t$. With $R_t=1000+50t$,
\[
\sum_{t=1}^{12}(1000+50t)=12\cdot1000+50\sum_{t=1}^{12}t=12000+50\cdot\frac{12\cdot13}{2}=12000+3900=15900,
\]
and the average is $15900/12=1325$.
\end{example}

\begin{example}[Discounting a stream of payments]
An asset pays $\pounds100$ at the end of each of the next three years, and the interest rate is $r=5\%$. Its present
value is $\sum_{t=1}^{3}\frac{100}{(1.05)^t}$. Evaluate it to the nearest penny.
\solution
$\frac{100}{1.05}+\frac{100}{1.05^2}+\frac{100}{1.05^3}=95.238+90.703+86.384=272.32$ (to the nearest penny). A
closed form for sums of this geometric type is proved by induction in \cref{ch:ind}.
\end{example}

% --------------------------------------------------------------------
\section{Chapter review}
% --------------------------------------------------------------------
\subsection*{Summary}
$\sum_{j=k}^n a_j$ denotes $a_k+\dots+a_n$; the index is a dummy variable and the sum has $n-k+1$ terms. A sum with
$n=k$ has one term; a sum with $n=k-1$ has none and is defined to be $0$, the only value compatible with P3. A sum of
a constant is the number of terms times the constant. P1 and P2 --- sums split over addition and constants come
outside --- follow from the properties of the real numbers; P3 --- adjacent ranges can be joined --- follows from the
definition. Products and squares cannot be taken through a sum.

\subsection*{Key definitions}
\begin{itemize}
  \item $\sum_{j=k}^{n}a_j=a_k+a_{k+1}+\dots+a_n$; index $j$, limits $k$ and $n$.
  \item $\sum_{j=k}^{k}a_j=a_k$; \ $\sum_{j=k}^{k-1}a_j=0$ (empty sum, by definition).
\end{itemize}

\subsection*{Key equations}
\begin{gather*}
\sum_{j=k}^{n}a=(n-k+1)a,\qquad \sum(a_j+b_j)=\sum a_j+\sum b_j,\qquad \sum c\,a_j=c\sum a_j,\\
\sum_{j=k}^{n}a_j+\sum_{j=n+1}^{m}a_j=\sum_{j=k}^{m}a_j\ \ (k\le n\le m).
\end{gather*}

\subsection*{Connections}
Sums are the natural setting for induction (\cref{ch:ind}): a statement about $\sum_{j=1}^{n}$ for every $n$ is proved
by peeling off the last term with P3. The binomial formula (\cref{ch:binom}) and polynomials (\cref{ch:func}) are
written with $\sum$. Infinite sums, the subject of Lecture~3, are limits (\cref{ch:limits}) of finite sums.

\subsection*{Common mistakes}
\begin{itemize}
  \item Counting $n-k$ terms instead of $n-k+1$.
  \item Starting at $j=0$ or $j=1$ when the lower limit is negative, and losing terms.
  \item Taking a factor depending on $j$ outside the sum.
  \item Writing $\sum a_jb_j=\sum a_j\sum b_j$ or $\sum a_j^2=(\sum a_j)^2$.
\end{itemize}

\subsection*{Exam checklist}
\begin{checklist}
  \item Expand and evaluate a sum with any integer limits.
  \item State the one-term and empty-sum conventions and justify the empty-sum value.
  \item State P1, P2 and P3 and say which come from the properties of $\R$ and which from the definition.
  \item Use the properties and the constant-sum rule to find closed forms.
\end{checklist}

% --------------------------------------------------------------------
\section{Practice questions}
% --------------------------------------------------------------------
\pq[basic] Evaluate $\sum_{j=1}^{5}(3j-2)$, $\sum_{j=0}^{4}2^j$ and $\sum_{j=-2}^{2}(j^2+1)$.

\pq[basic] Write in summation notation: (a) $1+4+9+16+25+36$; (b) $2+4+6+\dots+100$; (c) $1-1+1-1+1-1$.

\pq[conceptual] How many terms does $\sum_{j=-3}^{10}a_j$ have? Explain why the formula for the number of terms
has a ``$+1$''.

\pq[mathematical] Using $\sum_{j=1}^n j=\frac{n(n+1)}{2}$ and the properties of $\sum$, find closed forms for
$\sum_{j=1}^{n}(6j-1)$ and $\sum_{j=11}^{20}j$.

\pq[mathematical] Show that $\sum_{j=1}^{n}\big(a_{j+1}-a_j\big)=a_{n+1}-a_1$ for any numbers $a_1,\dots,a_{n+1}$.
Use this to evaluate $\sum_{j=1}^{n}\Big(\frac{1}{j}-\frac{1}{j+1}\Big)$.

\pq[application] A firm's cost of producing the $j$-th unit is $c_j=20+0.5j$ (so the first unit costs $20.5$). Write
the total cost of producing $n$ units as a sum and find a closed form. How much does it cost to produce $40$ units?

\pq[multi-step] Show by a counterexample that $\sum_{j=1}^{n}a_jb_j=\Big(\sum_{j=1}^{n}a_j\Big)\Big(\sum_{j=1}^{n}b_j\Big)$
is false in general, and find a condition on $n$ under which it is always true.

\pq[exam-style] (a) State the value of $\sum_{j=7}^{6}a_j$ and explain why this value is adopted.
(b) Suppose $a_j=c$ for every $j$. Explain why $\sum_{j=k}^{n}a_j=(n-k+1)c$ and evaluate $\sum_{j=-5}^{9}4$.
(c) Find a closed form for $\sum_{j=1}^{n}(4j+3)$, naming each property of the summation operator used.

% --------------------------------------------------------------------
\section{Solutions to practice questions}
% --------------------------------------------------------------------
\pqsol{6.1}
$1+4+7+10+13=35$; \ $1+2+4+8+16=31$; \ $5+2+1+2+5=15$.

\pqsol{6.2}
(a) $\sum_{j=1}^{6}j^2$. (b) $\sum_{j=1}^{50}2j$. (c) $\sum_{j=1}^{6}(-1)^{j+1}$ (or $\sum_{j=0}^{5}(-1)^j$).

\pqsol{6.3}
$10-(-3)+1=14$ terms. The indices run $-3,-2,\dots,10$. The difference $n-k=13$ counts the steps between the first
and last index; the number of indices is one more than the number of steps, because both endpoints are included.

\pqsol{6.4}
$\sum_{j=1}^{n}(6j-1)=6\cdot\frac{n(n+1)}2-n=3n^2+3n-n=3n^2+2n$.
By P3, $\sum_{j=11}^{20}j=\sum_{j=1}^{20}j-\sum_{j=1}^{10}j=210-55=155$.

\pqsol{6.5}
Write out the sum: $(a_2-a_1)+(a_3-a_2)+\dots+(a_{n+1}-a_n)$. Each $a_j$ with $2\le j\le n$ appears once with $+$ and
once with $-$, so everything cancels except $a_{n+1}-a_1$. (Formally: by P1 and P2 the sum is
$\sum_{j=1}^n a_{j+1}-\sum_{j=1}^n a_j=\sum_{i=2}^{n+1}a_i-\sum_{j=1}^n a_j=a_{n+1}-a_1$, using an index shift and
P3.) With $a_j=-\frac1j$: $\frac1j-\frac1{j+1}=a_{j+1}-a_j$, so the sum is $a_{n+1}-a_1=1-\frac1{n+1}=\frac{n}{n+1}$.

\pqsol{6.6}
$C(n)=\sum_{j=1}^{n}(20+0.5j)=20n+0.5\cdot\frac{n(n+1)}2=20n+\frac{n(n+1)}{4}$.
$C(40)=800+\frac{40\cdot41}{4}=800+410=1210$.

\pqsol{6.7}
$n=2$, $a_1=a_2=b_1=b_2=1$: the left side is $2$, the right side is $2\cdot2=4$. When $n=1$ both sides equal $a_1b_1$, so
the identity holds for every choice of numbers exactly when $n=1$ (for $n\ge2$ the example above, padded with zeros,
is a counterexample).

\pqsol{6.8}
(a) $0$. The upper limit is one less than the lower limit, so the sum has no terms, and a sum of zero terms is
defined to be $0$. The choice keeps P3 valid in the boundary case (\cref{prop:empty-forced}), and zero is the
number whose addition changes nothing.
(b) The sum is $c+c+\dots+c$ with one term for each index $k,k+1,\dots,n$, i.e.\ $n-k+1$ terms, so it equals $(n-k+1)c$.
Here $9-(-5)+1=15$ terms, so the sum is $60$.
(c) See \cref{ex:closed-form}: $2n^2+5n$.
